\section{Introduction}
\label{sec:intro}

Operational data assimilation (DA) is built on commitments so standard that they
are seldom stated as hypotheses at all. The analysis is produced by cycling:
propagate a state estimate, update it when observations arrive, repeat. Skill is
understood to follow model fidelity, so model development and DA development are
the same programme. The state is a point in the phase space of a differential
equation, and its uncertainty is a perturbation of that point --- a covariance,
or an ensemble. Prior errors travel by the same integration that carries the
signal.

Each of these buys something real, and none is foolish. Together, we argue, they
also explain two symptoms that the field observes repeatedly and treats as
tuning problems:

\begin{description}[leftmargin=2em,style=nextline]
  \item[\textnormal{\textbf{S-a --- the marginal value of observations.}}]
    DA converts additional sparse observations into skill inefficiently, and the
    inefficiency worsens sharply under model error. The quantity of interest is
    not RMSE at a given observation density but
    $\mathrm{d}\,\text{skill}/\mathrm{d}\,\text{density}$, measured at matched
    everything.
  \item[\textnormal{\textbf{S-b --- posterior quality.}}]
    DA's posterior characterisation is structurally impoverished, not merely
    mis-calibrated: spread/skill ratios miss their reliability target in ways
    that no inflation parameter repairs, because the deficit is a property of
    the estimator class.
\end{description}

\paragraph{What this paper is, and is not.}
This is a \emph{diagnosis} paper. We propose no new assimilation method. That is
a deliberate choice rather than a limitation of scope: a method paper invites
the reply ``your baseline was under-tuned'', whereas the claim here concerns
what the standard baselines cannot do \emph{however well tuned}. We therefore
report the tuning-sensitivity studies (Section~\ref{sec:tuning}) as evidence
that the baselines were mapped rather than assumed, including one that returned
a negative result.

\paragraph{Contributions.}
\begin{enumerate}[leftmargin=2em]
  \item We state four structural hypotheses (\textbf{H1}--\textbf{H3b},
    Section~\ref{sec:hypotheses}) in a form that admits falsification, and show
    that each can be relaxed \emph{separately} so that the resulting change in
    skill can be attributed to a named commitment.
  \item We show that model error \emph{collapses the marginal value of
    observations} rather than adding a constant error offset
    (Section~\ref{sec:marginal}): a $6.2\times$ collapse for strong-constraint
    4D-Var against $1.9\times$ for ensemble filters. To the best of our survey
    this effect is unreported.
  \item We show that the ODE representation's sensitivity to model parameters is
    \emph{not} a proxy for the posterior's sensitivity to the same parameters
    (Section~\ref{sec:sensitivity}), and that classical DA conflates the two by
    construction because its estimator \emph{is} the model. The consequence is
    methodological: assessing parameter uncertainty by propagating it through
    the model systematically overstates its effect on the posterior whenever the
    observations are informative.
  \item We show that the choice of state variable, at bit-identical physics,
    changes analysis skill substantially and predictably
    (Section~\ref{sec:representation}).
  \item We record a corollary about gradient-based assimilation
    (Section~\ref{sec:adjoint}): the value of the adjoint is conditional on
    model fidelity. At matched conditions the adjoint method and the best
    ensemble filter convert extra observations into skill at an identical rate;
    under model error the adjoint method's return collapses $3.3\times$ harder.
\end{enumerate}

\paragraph{Testbeds.}
All figures are from this study's own runs on two systems: a two-scale
Lorenz-96 model and a two-layer quasi-geostrophic (QG) channel. Both are
idealised. We take that as the price of being able to define ``matched
everything'' precisely enough for a marginal-value measurement to mean
something, and we return to the question of transfer in
Section~\ref{sec:limitations}.
